\chapter{Curvature Pinching and Preserved Curvature Properties}

\section{Curvature pinching and preserved curvature properties under Ricci flow}

The chapter develops a maximum-principle framework for curvature tensors
evolving by nonlinear heat equations. The main applications in dimension three
are the preservation of weakly positive sectional curvature, the preservation
of Ricci pinching inequalities, and a roundness estimate when scalar curvature
is large.

\subsection{The Einstein tensor and its evolution}

\begin{definition}
\label{topping-ch9-einstein-tensor}
\dcref{9.2.1}
\group{curvature-pinching}
\source{topping:page-0106,topping:page-0107}
The \emph{Einstein tensor} is the symmetric \((0,2)\)-tensor
\[
E := -\Ric + \frac{R}{2}\,g.
\]
In dimension three, if \(\lambda_1,\lambda_2,\lambda_3\) are the sectional
curvatures of an orthonormal basis of \(2\)-planes, then \(E\) is diagonal in
the corresponding orthonormal basis and its eigenvalues are
\(\lambda_1,\lambda_2,\lambda_3\).
\end{definition}

\begin{proposition}
\label{topping-ch9-einstein-evolution}
\dcref{9.3.1,9.3.2}
\group{curvature-pinching}
\source{topping:page-0107}
\uses{topping-ch9-einstein-tensor}
Under Ricci flow, the Einstein tensor satisfies
\[
\left\langle \frac{\partial E}{\partial t}(X),W\right\rangle
= \Delta E(X,W)-2\langle \Rm(X,\cdot,W,\cdot),\Ric\rangle + |\Ric|^2 g(X,W),
\]
and, in an orthonormal frame diagonalising the curvature operator,
\[
\frac{\partial E}{\partial t}-\Delta E
= 2\begin{pmatrix}
\lambda_1^2+\lambda_2\lambda_3 & 0 & 0\\
0 & \lambda_2^2+\lambda_1\lambda_3 & 0\\
0 & 0 & \lambda_3^2+\lambda_1\lambda_2
\end{pmatrix}.
\]
\end{proposition}

\subsection{The Uhlenbeck trick}

\begin{definition}
\label{topping-ch9-uhlenbeck-setup}
\dcref{9.4.1}
\group{curvature-pinching}
\source{topping:page-0108,topping:page-0109}
Let \((\mathcal M,g(t))\) be a Ricci flow and let \(V\to\mathcal M\) be an
abstract vector bundle with a bundle isomorphism \(u_0:V\to T\mathcal M\)
compatible with the initial metric. The time-dependent bundle isomorphisms
\(u_t:V\to T\mathcal M\) are defined by
\[
\frac{\partial u}{\partial t}=\Ric_{g(t)}(u),\qquad u(0)=u_0.
\]
The pull-back connection \(A(t)\) on \(V\) is given by
\[
u(A_X v)=\nabla_X(u(v)).
\]
\end{definition}

\begin{remark}
\label{topping-ch9-uhlenbeck-metric}
\dcref{9.4.2,9.4.3}
\group{curvature-pinching}
\source{topping:page-0108,topping:page-0109}
The family \(u_t\) remains a fibrewise isometry, so the induced metric \(h\)
on \(V\) is independent of time and the connection \(A(t)\) is compatible with
\(h\). Moreover, \(Au=0\) and \(Ah=0\).
\end{remark}

\begin{definition}
\label{topping-ch9-pulled-back-einstein}
\dcref{9.4.4,9.4.5}
\group{curvature-pinching}
\source{topping:page-0109,topping:page-0110}
Let \(\widetilde E\in\Gamma(V^*\otimes V)\) be the pull-back of the Einstein
tensor under \(u\), defined by
\[
u(\widetilde E(v))=E(u(v)).
\]
Then \(\widetilde E\) is a section of the bundle of \(h\)-symmetric endomorphisms
and satisfies the same evolution equation as \(E\) with \(A(t)\) in place of
\(\nabla\).
\end{definition}

\subsection{Parallel functions on vector bundles}

\begin{definition}
\label{topping-ch9-parallel-function}
\dcref{9.5.2}
\group{curvature-pinching}
\source{topping:page-0111}
Let \(W\to\mathcal M\) be a vector bundle with fibre metric \(h\) and
compatible connection \(A(t)\). A smooth function \(\Psi:W\to\mathbb R\) is
\emph{parallel} if it is invariant under parallel translation with respect to
\(A(t)\) at each time.
\end{definition}

\begin{definition}
\label{topping-ch9-convex-parallel-function}
\dcref{9.5.3}
\group{curvature-pinching}
\source{topping:page-0112}
A parallel function \(\Psi:W\to\mathbb R\) is called \emph{convex} if its
restriction to any fibre \(W_p\) is convex.
\end{definition}

\begin{proposition}
\label{topping-ch9-parallel-composition-parabolic}
\dcref{9.5.4}
\group{curvature-pinching}
\source{topping:page-0112,topping:page-0113}
\uses{topping-ch9-parallel-function,topping-ch9-convex-parallel-function}
If \(E(t)\in\Gamma(W)\) satisfies
\[
\frac{\partial E}{\partial t}-\Delta_AE=\Upsilon(E),
\]
and \(\Psi\) is parallel and convex, then
\[
\left(\frac{\partial}{\partial t}-\Delta_{\mathcal M}\right)(\Psi\circ E)
\le d\Psi(E)\!\left(\Upsilon(E)\right).
\]
\end{proposition}

\begin{definition}
\label{topping-ch9-parallel-convex-subset}
\dcref{9.5.5}
\group{curvature-pinching}
\source{topping:page-0113}
A subset \(X\subset W\) is \emph{parallel} if it is invariant under parallel
translation, and \emph{convex} if its intersection with each fibre is convex.
\end{definition}

\subsection{The ODE-PDE theorem}

\begin{theorem}
\label{topping-ch9-ode-pde-theorem}
\dcref{9.6.1}
\group{curvature-pinching}
\source{topping:page-0114,topping:page-0115,topping:page-0116}
\uses{topping-ch9-parallel-composition-parabolic,topping-ch9-parallel-convex-subset}
Let \(W\to(\mathcal M,g(t))\) be a bundle with compatible connection \(A(t)\),
and let \(E(t)\in\Gamma(W)\) satisfy
\[
\frac{\partial E}{\partial t}-\Delta_AE=\Upsilon(E).
\]
If \(X\subset W\) is closed, convex, and parallel, and the associated fibrewise
ODE
\[
\frac{de}{dt}=\Upsilon(e)
\]
preserves \(X\), then \(E(\cdot,0)\subset X\) implies \(E(\cdot,t)\subset X\) for
all \(t\in[0,T]\).
\end{theorem}

\begin{remark}
\label{topping-ch9-ode-pde-boundary-test}
\dcref{9.6.2}
\group{curvature-pinching}
\source{topping:page-0116}
For a set \(X\) given as the \(0\)-sublevel set of a convex parallel function
\(\Psi\), it suffices to check the boundary derivative condition
\[
\left.\frac{d}{dt}\Psi(e(t))\right|_{t=0}\le 0
\]
for the ODE solutions starting at points of \(\partial X\).
\end{remark}

\subsection{Applications to curvature pinching}

\begin{remark}
\label{topping-ch9-ode-eigenvalue-system}
\dcref{9.7.1}
\group{curvature-pinching}
\source{topping:page-0116,topping:page-0117}
For the pulled-back Einstein tensor, the fibrewise ODE preserves eigenvectors
and reduces to
\[
\frac{d\lambda_1}{dt}=2(\lambda_1^2+\lambda_2\lambda_3),\quad
\frac{d\lambda_2}{dt}=2(\lambda_2^2+\lambda_3\lambda_1),\quad
\frac{d\lambda_3}{dt}=2(\lambda_3^2+\lambda_1\lambda_2).
\]
\end{remark}

\begin{lemma}
\label{topping-ch9-eigenvalue-convexity}
\dcref{9.7.2}
\group{curvature-pinching}
\source{topping:page-0117,topping:page-0118}
On the vector space of symmetric endomorphisms of \(\mathbb R^3\), the trace
\(\lambda_1+\lambda_2+\lambda_3\) is linear and the least eigenvalue
\(\lambda_1\) is concave. Consequently, \(\lambda_3\) is convex,
\(\lambda_1+\lambda_2\) is concave, and \(\lambda_3-\lambda_1\) is convex.
If \(\alpha\) is concave and \(\beta\) is concave and increasing, then
\(\beta\circ\alpha\) is concave on the domain where it is defined.
\end{lemma}

\begin{theorem}
\label{topping-ch9-positive-sectional-preserved}
\dcref{9.7.3}
\group{curvature-pinching}
\source{topping:page-0118}
\uses{topping-ch9-eigenvalue-convexity,topping-ch9-ode-pde-theorem,topping-ch9-ode-eigenvalue-system}
Weakly positive sectional curvature is preserved under Ricci flow on
three-dimensional closed manifolds.
\end{theorem}

\begin{remark}
\label{topping-ch9-positive-sectional-strong}
\dcref{9.7.4}
\group{curvature-pinching}
\source{topping:page-0118}
A minor adaptation of the argument also preserves strictly positive sectional
curvature in the three-dimensional closed case.
\end{remark}

\begin{theorem}
\label{topping-ch9-ricci-pinching-preserved}
\dcref{9.7.5}
\group{curvature-pinching}
\source{topping:page-0119,topping:page-0120}
\uses{topping-ch9-eigenvalue-convexity,topping-ch9-ode-pde-theorem,topping-ch9-ode-eigenvalue-system}
For every \(\varepsilon\in[0,\tfrac13)\), the pinching condition
\[
\Ric\ge \varepsilon Rg
\]
is preserved under Ricci flow on three-dimensional closed manifolds.
\end{theorem}

\begin{corollary}
\label{topping-ch9-ricci-positive-preserved}
\dcref{9.7.6,9.7.7}
\group{curvature-pinching}
\source{topping:page-0120}
\uses{topping-ch9-ricci-pinching-preserved}
Under Ricci flow on a three-dimensional closed manifold, the conditions
\(\Ric> \varepsilon Rg\) for \(\varepsilon\in[0,\tfrac13)\), as well as
\(\Ric\ge 0\) and \(\Ric>0\), are preserved.
\end{corollary}

\begin{theorem}
\label{topping-ch9-roundness-estimate}
\dcref{9.7.8}
\group{curvature-pinching}
\source{topping:page-0121,topping:page-0122,topping:page-0123}
\uses{topping-ch9-eigenvalue-convexity,topping-ch9-ode-pde-theorem,topping-ch9-ode-eigenvalue-system,topping-ch9-ricci-pinching-preserved}
For any \(0<\beta<B<\infty\) and \(\gamma>0\), there exists
\(M=M(\beta,B,\gamma)<\infty\) such that if \(g(t)\) is a Ricci flow on a
closed three-dimensional manifold with
\(\beta g(0)\le \Ric(g(0))\le Bg(0)\), then
\[
\left|\Ric-\frac13 Rg\right|\le \gamma R+M
\]
for all \(t\in[0,T]\).
\end{theorem}

\begin{remark}
\label{topping-ch9-roundness-reduction}
\dcref{9.7.9}
\group{curvature-pinching}
\source{topping:page-0121}
It suffices to control the weakly positive quantity \(\lambda_3-\lambda_1\),
since
\[
\left|\Ric-\frac13Rg\right|^2\le (\lambda_3-\lambda_1)^2.
\]
\end{remark}
